\section{A spatially varying quiet fiber profile}\label{sec:quiet}

We now choose a radial fiber profile that tolerates spatial variation
of its parameters. Its steep part gives the curvature margins needed
near the testing disks; its logarithmic tail keeps the whole affine
fiber available. After establishing the scalar estimates, we use the
six-block formulas to prove a uniform curvature criterion and then
show that it survives localized shear terms.

Fix \(b_0\ge1000\) and \(K_0\ge1\); increasing \(K_0\) later is
allowed.  Define a positive function \(u\) of \(\tau>0\) by
\begin{equation}\label{quiet:u}
 u=\tau\exp\left(\int_0^\tau
       \frac{(1+v)^{-1/16}-1}{v}\,dv\right).
\end{equation}
It satisfies \(u/\tau\to1\) at zero and
\(d\log u/d\log\tau=(1+\tau)^{-1/16}\).  It is increasing, is
comparable to \(\tau\) on \((0,1]\), and tends to a finite positive
limit at infinity: the last assertion follows by integrating
\((1+\tau)^{-1/16}d\tau/\tau\) from 1 to infinity.
More specifically, convexity gives
\((1+v)^{-1/16}\ge1-v/16\), whence
\(e^{-1/16}\tau\le u\le\tau\) for \(0<\tau\le1\).
For a real parameter \(E\), called the severity, set
\begin{equation}\label{quiet:q}
 q=q_0(x,E)=(1+e^E u)^{b_0}+K_0^{-1}\log(1+\tau),
 \qquad x=\log\tau.
\end{equation}
Constants in this section may depend on \(b_0,K_0\), but are independent
of the spatial dimension.  A lower threshold for \(E\) will always be
chosen sufficiently large for these constants.  The primitive \(F\) has
the normalization \eqref{fiber:F}.  Because \(q-1=O_E(\tau)\) at zero,
this normalization exists.  Because \(q\) grows only linearly in \(x\) at
the right end, \(F\to+\infty\) there.

\begin{lemma}[Scalar estimates for the profile]\label{quiet:scalar}
Put
\[
 q_a=(1+e^Eu)^{b_0},\quad r=r_q=q_x/q,\quad
 D=D_q=(q_x+q_{xx})/q=r(1+r)+r_x,\quad e_0=q_E/q.
\]
For all sufficiently large \(E\) and all \(\tau>0\),
\begin{gather}
 q\ge1,\quad q_x>0,\quad
 D\ge r(r+\tfrac12),\quad q_{xx}/q_x\ge-\tfrac1{16},
 \label{quiet:slopes}\\
 0<r\le C,\quad |e_0|+|q_{EE}/q|\le C,
 \quad |r_E|\le Cr.
 \label{quiet:derivatives}
\end{gather}
For \(\tau\le1\), the first two quantities in the middle of
\eqref{quiet:derivatives} are also at most \(Cr\).
For \(1\le j\le4\),
\begin{equation}\label{quiet:Fderivatives}
 |\partial_E^jF|+|\partial_E^j\bar F|\le C,
 \qquad \bar F(\tau,E)=\tau^{-1}\int_0^\tau F(\log v,E)\,dv.
\end{equation}
For \(0\le j\le4\) and \(\tau\ge e^{-E/4}\),
\begin{equation}\label{quiet:difference}
 |\partial_E^j(F-\bar F)|\le C e^{-E/4},\qquad
 |\partial_E^j(F-\bar F)|\le Cr.
\end{equation}
The following additional estimates will be used:
\begin{align}
 q_x&\ge e^E &&(e^{-E/4}\le\tau\le e^E),
 \label{quiet:intermediate}\\
 q&\le C(e^{b_0E}+\log(1+\tau)),\qquad
 q_x\ge c(1+e^{b_0E}(1+\tau)^{-1/16}) &&(\tau\ge1),
 \label{quiet:tail}\\
 r&\ge .9 b_0(1+\tau)^{-1/16}
              &&(1\le\tau\le e^{3b_0E}),
 \label{quiet:middle-slope}\\
 q&\le C,\quad q-1\le Ce^E\tau,\quad
 r\asymp e^E\tau,\quad D\asymp e^E\tau
              &&(0<\tau\le2e^{-E}).
 \label{quiet:axis-slope}
\end{align}
In particular \(q_x\ge c\) for \(\tau\ge e^{-E}\), and
\(q^{-1}\le Cr\) for \(\tau\ge2e^{-E}\).
\end{lemma}

\begin{proof}
Write \(a=(1+\tau)^{-1/16}\), \(z=e^Eu\), and
\(q_b=K_0^{-1}\log(1+\tau)\).  The logarithmic slopes of the two
summands of \(q\) are
\[
 p=\frac{(q_a)_x}{q_a}=b_0a\frac z{1+z},\qquad
 d=\frac{(q_b)_x}{q_b}
     =\frac{\tau}{(1+\tau)\log(1+\tau)}.
\]
They satisfy
\[
 p_x/p\ge-\tfrac1{16},\qquad
 d_x/d=\frac1{1+\tau}
       -\frac{\tau}{(1+\tau)\log(1+\tau)}\ge-\tfrac12.
\]
For the second inequality use
\(\log(1+\tau)\ge2\tau/(2+\tau)\), obtained by differentiation from
equality at zero.  The derivative of a weighted mean of logarithmic
slopes is the weighted mean of their derivatives plus their nonnegative
variance.  Hence \(r_x\ge-r/2\), proving the first bound for \(D\).
Also
\((q_a)_{xx}/(q_a)_x=p+p_x/p\ge-1/16\) and
\((q_b)_{xx}/(q_b)_x=(1+\tau)^{-1}>0\).  Taking the weighted mean with
weights \((q_a)_x,(q_b)_x\) gives the other slope bound.

Every positive \(E\) derivative of \(q_a\) of order at most four is bounded
by \(Cq_a\min(1,z)\).  Differentiating \(q^{-1}\) thus gives
\begin{equation}\label{quiet:reciprocal}
 |\partial_E^j(q^{-1})|
       \le C\frac{q_a}{q^2}\min(1,z),\qquad 1\le j\le4.
\end{equation}
Likewise \((q_a)_x\) and its \(E\) derivative have ratio bounded in
absolute value by a profile constant, so \(|r_E|\le Cr\).
Since \(a\ge2^{-1/16}\) for \(\tau\le1\), comparison with \(p\) gives
\(e_0,q_{EE}/q\le Cr\) there.  The remaining bounds in
\eqref{quiet:derivatives} follow directly from the displayed formulas.

Differentiate the normalized integral for \(F\).  For \(x\le0\), use
\(d x=d\log z/a\), where \(a\ge2^{-1/16}\), in
\eqref{quiet:reciprocal}; the resulting absolute integral is at most
\[
 C\int_0^\infty\frac{\min(1,z)}{(1+z)^{b_0}}\,\frac{dz}{z}<\infty.
\]
For \(x\ge0\), put \(M=e^{b_0E}\).  The positive lower and upper
bounds on \(u\) for \(\tau\ge1\) imply
\(cM\le q_a\le CM\) and
\(q\ge cM+x/K_0\).  The bound in \eqref{quiet:reciprocal} is therefore
at most \(CM/(cM+x/K_0)^2\), whose integral over \([0,\infty)\)
is uniformly bounded.  This proves \eqref{quiet:Fderivatives}, also
for the average by integration.

Integration by parts gives, including for \(j=0\),
\begin{equation}\label{quiet:averageidentity}
 \partial_E^j(F-\bar F)
   =\int_0^1\partial_E^j(q^{-1})(\log\tau+\log t,E)\,dt.
\end{equation}
If \(\tau\ge e^{-E/4}\), split at \(t=e^{-E/4}\).  On the first
part the integrand is uniformly bounded.  On the second part
\(\tau t\ge e^{-E/2}\), and
\(u(v)\ge c\min(v,1)\) implies \(e^E u(\tau t)\ge c e^{E/2}\).
There the integrand, also for \(j=0\), is at most \(C e^{-b_0E/2}\).
This proves the first bound in \eqref{quiet:difference}.
For \(\tau\ge2\), split instead at \(t=\tau^{-1/2}\).
Monotonicity and the positive finite limit of \(u\), together with
\(\log(1+\sqrt\tau)\ge\tfrac12\log(1+\tau)\), give
\(q(\log\sqrt\tau,E)\ge c q(\log\tau,E)\).  Thus the integral is
at most \(C(\tau^{-1/2}+q^{-1})\).
The estimates
\(q\le C(M+\log\tau)\) and
\(q_x\ge c(1+M\tau^{-1/16})\) show that
\(q\tau^{-1/2}\le Cq_x\) and \(1\le Cq_x\); hence this last bound is
at most \(Cr\).  On \([e^{-E/4},2]\) the explicit formula for \(p\) gives a
positive lower bound on \(r\).  This completes \eqref{quiet:difference}.

The estimates \eqref{quiet:tail} follow from the same two summands of \(q\).
For \(1\le\tau\le e^{3b_0E}\), the logarithmic summand divided by
\(q_a\) tends uniformly to zero as the lower threshold on \(E\) increases;
also \(z/(1+z)\to1\) uniformly.  This proves
\eqref{quiet:middle-slope}.  On
\([e^{-E/4},e^E]\) one has
\(z\ge c e^{3E/4}\) and \(a\ge c e^{-E/16}\), so
\[
 q_x\ge(q_a)_x=b_0az(1+z)^{b_0-1}
     \ge c\exp((3b_0/4-1/16)E)\ge e^E.
\]
Finally, on \((0,2e^{-E}]\), the quantity \(z\) is comparable to
\(e^E\tau\) and remains bounded by a profile constant.  The formulas
for the first two \(x\) derivatives of \(q\) then give
\eqref{quiet:axis-slope}.  Between \(2e^{-E}\) and 1 the same formulas
give \(q_x\ge c\); for \(\tau\ge1\), use \eqref{quiet:tail}.
These statements also prove the two final consequences.
\end{proof}

The scalar estimates make two tasks compatible: the spatial Schur
metric remains positive, and the fiber curvature absorbs the trace
couplings in a real-plane test. The following criterion states the
spatial hypotheses needed for both tasks. The word \emph{quiet} refers
to the smallness of the severity derivatives in the metric of the twist.

\begin{proposition}[Quiet profile criterion]\label{quiet:profile}
There exist positive constants \(\epsilon,\mu\) and a lower threshold
\(E_{\min}\), depending on the fixed profile constants but not on \(n\),
with the following property.  Let \(E\) be a smooth spatial function and
use \(q=q_0(x,E(s))\) with constant center and no shear.  At a spatial
point suppose, in some fixed holomorphic affine coordinates, that:
\begin{enumerate}
\item \(G\ge V>0\).  The operator norms of \(V,G\), their inverses,
their first and second metric jets, and all derivatives of \(E\) of orders
one through four are at most \(e^{\epsilon E}\).
\item For every Hermitian contravariant matrix \(T\) of rank at most two,
\(C^V(T,T)\ge e^{-\epsilon E}|T|^2\), with the norm in those affine
coordinates.
\item In \(G\)-unit axes,
\[
 |\partial E|_G^2\le\mu,\qquad
 \|\partial\bar\partial E\|_G\le\mu,\qquad
 \|\nabla_h^{2,0}E\|_G\le\mu
 \quad\hbox{for every }\tau>0.
\]
\end{enumerate}
If \(E\ge E_{\min}\), the resulting metric is positive for every \(y\)
over that spatial point and has nonpositive real sectional curvature
there.  Uniformly in \(\tau\),
\begin{equation}\label{quiet:schur}
 \mathcal N=(1+O(\mu))G,\qquad
 h\asymp V+\tau G.
\end{equation}
More precisely, off the fiber axis the rank-two curvature calculation
retains fixed positive multiples of
\begin{equation}\label{quiet:margins}
 At^2+W(b_\times,\bar b_\times)+P|T|_{\mathcal N}^2
 +\boldsymbol{1}_{\{\tau\le e^{-E/4}\}}
                          e^{-\epsilon E}|T|^2.
\end{equation}
Here the fiber blocks use the tilted logarithmic coordinates,
\(|T|_{\mathcal N}\) is its Frobenius norm in \(\mathcal N\)-unit
spatial axes, and \(\boldsymbol{1}\) denotes the indicated indicator.
The meaning of the final hypothesis is noncircular: positivity of \(h\)
follows from the preceding hypotheses and does not use
\(\nabla_h^{2,0}E\).
\end{proposition}

\begin{proof}
Differentiating \(U_\tau=\psi-F\) gives
\begin{equation}\label{quiet:N}
 \mathcal N=G-F_E\partial\bar\partial E
               -F_{EE}\partial E\otimes\bar\partial E,
 \qquad h=V+\tau\bar{\mathcal N}.
\end{equation}
The bar denotes averaging from 0 to \(\tau\).  The bounded derivatives
in \eqref{quiet:Fderivatives} give
\((1-C\mu)G\le\mathcal N,\bar{\mathcal N}\le(1+C\mu)G\).
Choose \(\mu\) small enough that \(C\mu<1/10\).  This proves
positivity and \eqref{quiet:schur}.

We will repeatedly use finite-jet estimates of the following precise
kind.  Differentiating \eqref{quiet:N} at most twice spatially gives
sums with at most four derivatives of \(E\) and at most four \(E\) derivatives
of \(F\) or \(\bar F\).  Each summand is a product of a fixed number of
the assumed jets.  The formulas for the inverse, connection, and
curvature add a fixed number of inverse metric factors.  Consequently
there is an absolute finite constant \(C_*\), independent of dimension,
such that these estimates cost at most \(C e^{C_*\epsilon E}\), with
the factors \(1+\tau\) or \((1+\tau)^{-1}\) displayed explicitly
below.  This conclusion follows by the product rule and
\(h^{-1}\le C e^{\epsilon E}(1+\tau)^{-1}I\); it uses operator norms,
not sums of tensor entries.  The same observation applies after
normalizing spatial axes by \(\mathcal N\), because
\(e^{-\epsilon E}I/C\le\mathcal N\le C e^{\epsilon E}I\).
Enlarge \(C_*\) once to cover these finitely many products, and choose
\(\epsilon>0\) so small that \(C_*\epsilon<1/100\).
Any profile-dependent multiplicative constants are absorbed by raising
\(E_{\min}\).

In \(\mathcal N\)-unit axes define
\begin{equation}\label{quiet:Y}
 \mathbf Y=I-\tau\mathcal N^{1/2}h^{-1}\mathcal N^{1/2}.
\end{equation}
The displayed formula indicates the normalization; equivalently it is
\(I-\tau h^{-1}\) after \(\mathcal N=I\) has been imposed in the
fixed axes at the point.  We claim
\begin{equation}\label{quiet:Ybounds}
 \|\mathbf Y\|\le r/20\quad(\tau\ge2e^{-E}),\qquad
 \mathbf Y=I+O(e^{-(1-C_*\epsilon)E})
                         \quad(\tau\le2e^{-E}).
\end{equation}
Always \(\|\mathbf Y\|\le1+C\mu\), because
\(h\ge\tau\bar{\mathcal N}\) and
\(\bar{\mathcal N}\) is close to \(\mathcal N\).
On \([2e^{-E},1]\), the explicit slope \(p\) in the scalar lemma yields
\(r\ge .35b_0\) for all sufficiently large \(E\): indeed
\(u\ge e^{-1/16}\tau\), \(z\ge2e^{-1/16}\),
\(a\ge2^{-1/16}\), and
\(q_a/q\ge(1+\log2/K_0)^{-1}\ge(1+\log2)^{-1}\).
The product of these three lower bounds in
\(p(q_a/q)=b_0a[z/(1+z)](q_a/q)\) exceeds \(.35b_0\).
This proves the first bound on that interval.
For \(\tau\ge1\), \eqref{quiet:difference} gives
\(\|\bar{\mathcal N}-\mathcal N\|_{\mathcal N}\le C\mu r\).
Set \(h_0=V+\tau\mathcal N\).  The inverse identity and
\(h,h_0\ge c\tau\mathcal N\) show that replacing \(h\) by \(h_0\) in
\eqref{quiet:Y} costs at most \(C\mu r\).  In the normalized axes
\(V\le\gamma I\), where \(\gamma=(1-C\mu)^{-1}\), and hence
\[
 \|I-\tau h_0^{-1}\|\le\frac\gamma{\gamma+\tau}.
\]
This is at most \(r/40\) on \([1,e^{3b_0E}]\), by
\eqref{quiet:middle-slope} and \(b_0\ge1000\).
For \(\tau>e^{3b_0E}\), use \(r\ge c/q\) and
\(q/\tau\le C(e^{b_0E}+\log(1+\tau))/\tau\); the supremum of the
last expression on this range tends to zero with the lower threshold
on \(E\).  Thus the same estimate holds there.  Decreasing \(\mu\) pays
the inverse-identity error.  Finally, if \(\tau\le2e^{-E}\), then
\(h\ge V\) and the raw size bounds give
\(\tau\|\mathcal N^{1/2}h^{-1}\mathcal N^{1/2}\|
 \le C\tau e^{C_*\epsilon E}\), proving the second bound.

Since \(P_\tau=q(1+r)\) and
\(P^2P_{\tau\tau}=Pq^2D\), the curvature blocks in
Lemma~\ref{fiber:blocks} become
\begin{align}
 \frac W{Pq}&=rI+\mathbf Y+q^{-1}\partial\bar\partial\log q,
 \label{quiet:Wformula}\\
 P^{-1}\nabla_h^{2,0}P
     &=e_0\nabla_h^{2,0}E+(q_{EE}/q)\partial E\otimes\partial E.
 \label{quiet:spinformula}
\end{align}
The final term in \eqref{quiet:Wformula} and the spin form divided by
\(Pq\) have norm at most \(C\mu r\) when \(\tau\ge2e^{-E}\),
and at most \(C\mu\) on the remaining interval.  This follows from
\eqref{quiet:derivatives}, the three smallness assumptions on \(E\), and
\(q^{-1}\le Cr\) away from the exceptional interval.
After decreasing \(\mu\),
\begin{equation}\label{quiet:Wbounds}
 W\asymp Pqr\,\mathcal N\quad(\tau\ge2e^{-E}),\qquad
 W\asymp Pq(1+r)\,\mathcal N\quad(\tau\le2e^{-E}),
\end{equation}
and the spin inequality \eqref{fiber:spinabs} holds.

For the fiber diagonal term,
\[
 \frac{|\partial P|_h^2}{Pq^2D}
 \le C\mu\frac{e_0^2}{qD}\le C\mu.
\]
Here \(h\ge c\tau G\); the final scalar quotient is bounded using
\(e_0\le Cr\) for \(\tau\le1\), and
\(qD\ge cqr=cq_x\ge c\) for \(\tau\ge1\).
Thus
\begin{equation}\label{quiet:Abounds}
 A=(1+O(\mu))Pq^2D>0.
\end{equation}
In the same normalized axes the cancellation in
\eqref{fiber:blockAi} gives
\[
 |A_{\cdot}|_{\mathcal N}
 \le CPq\sqrt\mu\bigl(|e_0|(r+\|\mathbf Y\|)+|r_E|\bigr)
 \le CPq\sqrt\mu\,r.
\]
The last inequality uses \(e_0\le Cr\) on the exceptional interval.
Combining with \eqref{quiet:Wbounds} and
\eqref{quiet:Abounds}, the squared ratio to the dual norm determined
by \(AW\) is at most \(C\mu\).  Consequently
\eqref{fiber:Aiabs} holds with any fixed desired small \(\eta>0\)
after decreasing \(\mu\).

We next pay the trace coupling.  Outside the exceptional interval,
\eqref{quiet:Ybounds} and \eqref{quiet:Wformula} imply
\[
 |W:T|\le Pqr\left(|\operatorname{tr}_{\mathcal N}T|
                       +\frac{\sqrt2}{16}|T|_{\mathcal N}\right).
\]
By \(D\ge r^2\), its cost against .98A is at most
\begin{equation}\label{quiet:tracecost}
 \frac P{.98(1-C\mu)}
 \left(|\operatorname{tr}_{\mathcal N}T|
                       +\frac{\sqrt2}{16}|T|_{\mathcal N}\right)^2.
\end{equation}
For \(0\le s\le\sqrt2\),
\((s+\sqrt2/16)^2/(1+s^2)<.8\); this can also be checked by
differentiating, since its maximum on this interval occurs at the
right endpoint.  Because \(T\) has rank at most two, decreasing \(\mu\)
therefore makes \eqref{quiet:tracecost} at most
\(.85P((\operatorname{tr}_{\mathcal N}T)^2+|T|_{\mathcal N}^2)\).
These are exactly the positive symmetrized terms in
\eqref{fiber:blockL}.
On \(\tau\le2e^{-E}\), the bounds
\(D\ge c e^E\tau\), \(q,r\le C\), and
\eqref{quiet:Wbounds} instead give a trace cost at most
\(Ce^{-E}|T|_{\mathcal N}^2\).

To estimate the remaining \(C^h\), note from \eqref{quiet:N} that
the difference \(h-V=\tau\bar{\mathcal N}\) and its first two jets
are at most \(C\tau e^{C_*\epsilon E}\).
On \(\tau\le e^{-E/4}\), the curvature formula or
Lemma~\ref{fiber:perturbation} consequently gives
\[
 |C^h(T,T)-C^V(T,T)|
       \le C\tau e^{C_*\epsilon E}|T|^2
       \le\tfrac12e^{-\epsilon E}|T|^2.
\]
In the complementary range, the inverse bound and the jet bounds give
\[
 |C^h(T,T)|\le C(1+\tau)e^{C_*\epsilon E}|T|^2.
\]
This is arbitrarily small compared with \(P|T|_{\mathcal N}^2\):
indeed, for \(\tau\ge e^{-E/4}\),
\(\tau q_a/(1+\tau)\ge e^{2E}\) after increasing the threshold on \(E\),
by the explicit power \(b_0\ge1000\); normalization of \(T\) costs only
\(e^{C_*\epsilon E}\).  On the exceptional interval, its trace cost
\(Ce^{-E}|T|_{\mathcal N}^2\) is likewise negligible compared with
\(e^{-\epsilon E}|T|^2\).
We have thus obtained \eqref{fiber:reserve} with
\begin{equation}\label{quiet:Rreserve}
 \mathcal R(T)\ge cP|T|_{\mathcal N}^2
  +c\boldsymbol{1}_{\{\tau\le e^{-E/4}\}}
                          e^{-\epsilon E}|T|^2,
\end{equation}
before paying the \(K\) coupling.

For that coupling split \(K\) into \(P\nabla^h\mathcal N\) and the two
terms containing \(P_i=Pe_0E_i\).  The latter have normalized size at
most \(CP|e_0|\sqrt\mu\); their squared dual \(W\) norm after pairing
with \(T\) is at most \(C\mu P|T|_{\mathcal N}^2\).  To see this directly
outside the exceptional interval, the scalar quotient to bound is
\(e_0^2/(qr)\); it is bounded by \eqref{quiet:derivatives} and
\eqref{quiet:tail}.  On the exceptional interval replace \(r\) in that
denominator by \(1+r\), which is easier.
For \(P\nabla^h\mathcal N\) use three ranges:
\begin{enumerate}
\item If \(\tau\le e^{-E/4}\), then \(W\ge cPq\mathcal N\), and
the raw jet bounds give squared dual \(W\) cost at most
\(C\tau e^{C_*\epsilon E}|T|^2\), negligible compared with the
second term of \eqref{quiet:Rreserve}.
\item If \(e^{-E/4}\le\tau\le e^E\), use
\(\|\nabla^h\mathcal N\|_{\mathcal N}\le Ce^{C_*\epsilon E}\).
The squared dual cost is at most
\(CPe^{C_*\epsilon E}(qr)^{-1}|T|_{\mathcal N}^2\).
By \eqref{quiet:intermediate} this is negligible compared with
the first term of \eqref{quiet:Rreserve}.
\item If \(\tau>e^E\), the identity \(\nabla^h h=0\) permits
\[
 \nabla^h\mathcal N
   =\nabla^h(\mathcal N-\bar{\mathcal N}-V/\tau).
\]
Equation \eqref{quiet:difference}, differentiated spatially once,
and the raw bounds give a normalized norm at most
\(C e^{-E/4+C_*\epsilon E}\).  In making this estimate the connection
has raw norm at most \(Ce^{C_*\epsilon E}\), because \(h\) has first
jets at most \(C(1+\tau)e^{C_*\epsilon E}\) and inverse norm at most
\(Ce^{C_*\epsilon E}/(1+\tau)\).  Since \(qr=q_x\ge c\), the
squared dual \(W\) cost is again negligible compared with
\(P|T|_{\mathcal N}^2\).
\end{enumerate}
Fix \(\eta\le1/100\) in Lemma~\ref{fiber:ranktwo}.  Choose \(\mu\)
small enough for the terms proportional to \(\mu\), and then
\(E_{\min}\) large enough for the exponentially small terms, so that
\eqref{fiber:Kabs} holds.  That lemma proves positivity of every
rank-two contraction with the retained margins
\eqref{quiet:margins}, and hence nonpositive real sectional curvature.
Smoothness and positivity at the axis follow from
Lemma~\ref{fiber:metric}; curvature extends by continuity.
\end{proof}

To realize the disks of Proposition~\ref{grid:criterion}, we need to
add a term linear in the fiber variable whose coefficient varies over
the spatial base. A locally constant coefficient contributes only a
pluriharmonic potential. The issue is therefore its derivative support,
where the estimates below must hold for arbitrarily large fiber radius.

\begin{proposition}[Localized shear perturbations]\label{quiet:shear}
Under the hypotheses of Proposition~\ref{quiet:profile}, suppose \(E\) and
the center are constant on the spatial region under consideration.
Assume also that the first and pure second spatial derivatives of
\(\psi\), in the affine coordinates of that proposition, have norms
at most \(e^{\epsilon E}\).  Add
\[
 \ell_{\mathrm{sh}}=\operatorname{Re}(c_{\mathrm{sh}}(s)(y-\zeta)).
\]
On the support of its coefficient's derivatives suppose
\(\psi\ge-J\), where \(J\ge0\), and
\begin{equation}\label{quiet:shear-small}
 \|j^{\le4}c_{\mathrm{sh}}\|\le e^{-C E-J/2}.
\end{equation}
For a sufficiently large \(C\) depending only on the profile and the
constants in the quiet criterion, this perturbation preserves metric
positivity and nonpositive real sectional curvature for all \(y\).
Moreover its metric is at least one half the unperturbed metric, and
fixed fractions of the curvature margins remain.  The choice of \(C\) is
independent of dimension.
\end{proposition}

\begin{proof}
Where the coefficient is locally constant the added potential is
pluriharmonic, so it changes no metric.  It remains to estimate the
derivative collar in \eqref{quiet:shear-small}.  Since \(E\) is constant,
\(P_i=0\), \(\mathcal N=G\), and \(h=V+\tau G\).
The holomorphic tilt used in Lemma~\ref{fiber:blocks} has linear and
quadratic coefficients bounded by \(Ce^{\epsilon E}\), because
\(U_\tau=\psi-F\).  Its constant term may and will be zero.  Derivatives
through order four of the exponential of this fixed quadratic
polynomial at the point cost at most \(Ce^{C\epsilon E}\), by the
ordinary product rule.  These are fixed-order multilinear estimates.

First consider \(\tau\ge e^{-E}\) in tilted logarithmic coordinates.
The metric inverse has norm at most \(Ce^{C\epsilon E}/\tau\).
By \eqref{quiet:Abounds}, \eqref{quiet:Wbounds},
\eqref{quiet:Rreserve}, and \(q_x\ge c\) on this range, the retained
curvature on a rank-two matrix \(H\) is at least
\begin{equation}\label{quiet:large-margin}
 c\tau e^{-C_1E}|H|^2
\end{equation}
in these fixed coordinates, for a fixed \(C_1\).
The cubic table \eqref{fiber:cubics} shows explicitly that the only
possibly nonzero connection blocks are
\[
 \Gamma^0_{00}=P_\tau,\qquad
 \Gamma^j_{i0}=P(h^{-1}G)^j_i,\qquad
 \Gamma^k_{ij}=(\Gamma^h)^k_{ij}.
\]
Consequently their operator norm is at most
\(Ce^{C_1E}q\), since \(r\) is bounded.  Allowing a larger fixed power
of \(q\) in the following estimates would give the same result.

The inequality \(q\ge1\) and the left normalization imply \(F\le x\).
Thus
\[
 |y-\zeta|=e^{(F-\psi)/2}\le e^{-\psi/2}\sqrt\tau.
\]
After the tilt, all shear jets needed for the metric and curvature,
of orders two, three, and four, are therefore at most
\begin{equation}\label{quiet:shear-large-jets}
 C\sqrt\tau\exp((-C+C_2)E),
\end{equation}
where \(e^{-J/2}e^{-\psi/2}\le1\) was used.
For every fixed integer \(r_1\), the scalar estimate
\eqref{quiet:tail}, also valid as an upper bound for small \(\tau\),
gives
\begin{equation}\label{quiet:log-absorption}
 \sup_{\tau\ge e^{-E}}\frac{q^{r_1}}{\sqrt\tau}
       \le C_{r_1}e^{C_{r_1}E}.
\end{equation}
Indeed \(q^{r_1}\le C_{r_1}(e^{r_1b_0E}
+\log(1+\tau)^{r_1})\); on \([e^{-E},1]\) use
\(\tau^{-1/2}\le e^{E/2}\), and on \([1,\infty)\) the quotient
\(\log(1+\tau)^{r_1}/\sqrt\tau\) is bounded.

For clarity, this bound controls the quadratic perturbation as well
as the linear terms.  Write \(\delta=e^{-CE}\) and absorb all other
fixed \(E\) powers into \(e^{cE}\).  The relative metric error is bounded
by \(\delta e^{cE}/\sqrt\tau\).  Each linear curvature error from
\eqref{fiber:perturbbound}, divided by
\eqref{quiet:large-margin}, is bounded by
\(\delta e^{cE}q^{r_1}/\sqrt\tau\) for some fixed \(r_1\).
The quadratic contraction error before division by that margin is
at most \(\delta^2e^{cE}q^{2r_1}\); after division its bound is
\[
 \delta^2e^{c'E}\left(\frac{q^{r_1}}{\sqrt\tau}\right)^2.
\]
Equation \eqref{quiet:log-absorption} makes all these quantities
arbitrarily small by a single sufficiently large \(C\).  It follows that
the metric is at least half its old value and that the curvature
errors are smaller than a fixed fraction of the retained margin,
uniformly on this whole fiber range.

For \(0<\tau\le e^{-E}\), exponentiate the tilted logarithmic
coordinate and use the affine fiber coordinate with scale
\(e^{-(\psi(s_0)+E)/2}\) at the fixed spatial point \(s_0\).
Denote its value by \(z\).  The zero constant term of the tilt and
\eqref{quiet:axis-slope} give
\begin{equation}\label{quiet:small-z}
 |z|^2=e^{E+F}\asymp e^E\tau.
\end{equation}
For example, \(|F-x|\le Ce^E\tau\) follows by integration of
\((q^{-1}-1)d\tau/\tau\) on this range, and proves the comparison.
The metric blocks in affine coordinates are \(h\) and \(P/|z|^2\).
Both have inverse norm at most \(e^{C_3E}\).  The curvature tensor
transforms tensorially: its fiber diagonal block divides by \(|z|^4\),
its \(W\) block by \(|z|^2\), and its spatial block is unchanged.
Since
\[
 A\ge c\tau^2e^E,\qquad W\ge c\tau G,\qquad
 \mathcal R(T)\ge c e^{-\epsilon E}|T|^2,
\]
the retained rank-two curvature margin in affine coordinates is at
least \(e^{-C_3E}|H|^2\), after increasing \(C_3\).

The connection in the new coordinates is also bounded by
\(e^{C_3E}\), without an apparent singularity at \(z=0\).
Indeed the logarithmic block \(\Gamma^0_{00}=P_\tau\) becomes
\((P_\tau-1)/z\), and
\(P_\tau-1=O(e^E\tau)\); thus it is bounded by
\(C\sqrt{e^E\tau}\).  The mixed spatial block becomes
\(Ph^{-1}G/z\), bounded by \(e^{C_3E}\) using
\eqref{quiet:small-z}.  The spatial-output, two-fiber-input block is
zero because \(P_i=0\), and the spatial connection is already bounded.
These checks cover all cubic entries.  The same bounds extend to zero
by the smoothness proved in Lemma~\ref{fiber:metric}.

In these affine coordinates the shear jets of orders two through four
are bounded by
\[
 e^{-CE-J/2}e^{-(\psi(s_0)+E)/2}Ce^{C_4\epsilon E}
 \le C\exp(-(C+1/2-C_4\epsilon)E).
\]
The estimate includes derivatives of the tilt and uses that \(|z|\) is
bounded on this range.  Apply Lemma~\ref{fiber:perturbation}, with the
inverse metric, connection, and reciprocal margin bounded by fixed
powers of \(e^{C_3E}\).  Increasing the same \(C\) makes all relative
metric errors and all curvature errors smaller than one half their
respective retained bounds, including at the axis.  The estimates
overlap at \(\tau=e^{-E}\), so they cover the entire fiber.

Finally, all exponents just used arise from derivatives of order at
most four, a fixed finite number of products in
\eqref{fiber:perturbformula}, and operator norms of linear maps.
Rank-two contractions introduce only the absolute constants in
\eqref{fiber:ranknorm}.  None of these operations counts spatial
indices.  This proves the asserted independence from \(n\).
\end{proof}
